\bookchapter{第五十六章：琢磨與發佈——測試、調試與上架 Google Play}

\chapterintro{App 既成，繼之奈何？此章答「自能行」至「能交付」之三事：何以高效調試而得 bug、何以書測試以防回歸、何以送 App 上 Google Play 而使天下皆得下載之。上架非徒儀式之感——乃使一作化為一產之分水嶺也。}

\booksection{一、調試三板斧：Logcat、斷點、UI 檢查器}

\begin{center}
\begin{tabularx}{\linewidth}{|X|X|X|}
\hline
\textbf{工具} & \textbf{所為何事} & \textbf{何以用之} \tabularnewline
\hline
Logcat & 觀系統/自書之日誌 & 代碼中 Log.d("TAG","值=\$x")，面板搜 TAG 以濾 \tabularnewline
斷點調試 & 使程序停於某行，觀變量之值 & 行號旁一點而出紅點，Debug 模式而運，逐步行之 \tabularnewline
Layout Inspector & 觀運行中界面之構 & 菜單 Tools Layout Inspector，實時觀組件之樹 \tabularnewline
\hline
\end{tabularx}
\end{center}

\begin{center}\small\color{gray}調試三板斧\end{center}

\begin{tcbinfo}{小識：報錯何以讀？}
崩之後，Logcat 中有赤色堆棧（Stacktrace），自下而上覓首行以己包名起者（譬如 com.example.yourapp），彼行基本即案發之場。先觀「Exception 之類」（空指針乎？類型之轉乎？），復觀「at ... 汝之文件:行號」。此招能解崩潰之百分之八十。

\end{tcbinfo}

\booksection{二、測試：為代碼上保險}

手測不足恃——改此功能，彼者潛壞。測試者，「悔之藥」而兼「防彈之衣」也。安卓之測分三層，新手先精究其第一層：

\begin{center}
\begin{tabularx}{\linewidth}{|X|X|X|X|}
\hline
\textbf{層級} & \textbf{所測何物} & \textbf{置於何所} & \textbf{遲速} \tabularnewline
\hline
單元測試 & 純邏輯（計算、校驗） & app/src/test/ & 秒級 \tabularnewline
儀器測試 & 真機之上界面之行 & app/src/androidTest/ & 分鐘級 \tabularnewline
手動測試 & 真人之感、體驗 & —— & 恆須之 \tabularnewline
\hline
\end{tabularx}
\end{center}

\begin{center}\small\color{gray}測試三層\end{center}

\begin{lstlisting}[language=Kotlin]
// 將"金額字符串轉數字"此種邏輯抽成純函數
fun parseAmount(input: String): Double? {
    val v = input.trim().toDoubleOrNull()
    return if (v != null && v > 0) v else null
}

// 測試它（JUnit4）
@Test
fun emptyInputReturnsNull() {
    assertEquals(null, parseAmount(""))
}
@Test
fun negativeReturnsNull() {
    assertEquals(null, parseAmount("-5"))
}
@Test
fun normalValueParses() {
    assertEquals(12.5, parseAmount("12.5"))
}
\end{lstlisting}

記賬冊中金額之校驗、統計之算，皆當有單元測試罩之。行測試：右鍵 test 目錄 Run。綠而皆過，則可安然而改碼矣。

\booksection{三、簽名打包：為 App 蓋「印」}

安卓要求凡安裝之包必簽之——猶蓋公印，以證「此包吾所出，未遭改竄」。調試時 Android Studio 自用 debug 之簽；發佈須用己之 release 簽。其序：

\begin{enumerate}
  \item 菜單 Build Generate Signed Bundle / APK；
  \item 創 keystore 之文件（密鑰之庫）——設其密碼而妥藏之；
  \item 擇 release + AAB 之式，構建乃成。
\end{enumerate}

\textbf{避坑之警：keystore 失則 App 失}\par

上架之後更新 App，必用同一之簽，否則店將拒其新。keystore 之文件與密碼一失，汝惟易新應用而重上架——舊用戶皆無由接焉。務以 keystore 與密碼備於安隱之所（雲盤加密 + 記事之冊），此安卓開發者至昂之教訓之一也。

\booksection{四、上架 Google Play 之全程}

註冊開發者賬號

準備商店素材

上傳 AAB

填數據安全表

審覈（數小時\textasciitilde{}數天）

上架

\begin{center}\small\color{gray}圖 6：Google Play 上架六步\end{center}

\begin{itemize}
  \item 註冊：一次 25 美元之註冊費，須可用之 Google 賬號；
  \item 素材：應用之圖標（512×512）、截圖（至少二張）、功能之述、私隱政策之鏈（收數據者必填）；
  \item 數據安全之表：如實申報所收何數據——不連網之記賬冊可「零收集」；
  \item 審核：新賬號首番上架恆為人工所審，數時至數日不等，見駁則依議而改。
\end{itemize}

國內之用戶慎之：Google Play 須特殊網絡之境，且中國大陸之分發常行國內之店（華為、小米、應用寶之屬）。上架國內之店，其程相類，然各店自有其審之好——同一 AAB/APK 可分發於多店。以為練手，先通一家即勝矣。

\booksection{五、商店素材何由製之：第一印象即繫下載之數}

商店頁者，用戶遊店之時所見之櫥窗也。三材定其第一印象：

\begin{itemize}
  \item 圖標：512×512，勿以截圖充數。若無設計之能，可以 AI 之器生簡潔幾何之形，縮於手機桌面之上猶能一望而識者，方為合格；
  \item 截圖：至少二張，以四至六張為佳，以彰最核心之功能場景；每張配一淺白之語（「一秒記一筆」勝「高效記帳之方」十倍）；
  \item 描述：首句當明「此為何物、予誰用之」，勿堆「最專、最強」之詞——商店搜索之權重，主視關鍵詞與下載轉化，誠實勝於誇張。
\end{itemize}

\begin{tcbinfo}{小訣：以 A/B 之法觀同行}
上架之前，遍觀商店同類 App 之前十名，取其截圖比而較之：其圖標何風？其截圖售何賣點？其標題何書？模仿非抄襲也——學其套路而易以己之內容，此凡產品冷啟動至速之一課也。

\end{tcbinfo}

\booksection{五、性能優化與崩潰上報小清單}

App 上架之前，此二事宜各費半小時為之：

\booksubsection{一、性能自察四條}

\begin{enumerate}
  \item 列表當用 LazyColumn：長列切莫以 Column 強排之——懶載者，安卓性能之第一課也；
  \item 圖勿硬塞原圖：當用 Coil 圖庫自動縮放而緩存之，一張 5MB 之原圖足以令低階之機滯塞；
  \item 勿於主線程任重事：數據庫、網絡、大文件皆當於協程中為之；
  \item 當以 Profiler 量之：Android Studio 自帶之性能分析器能觀記憶體、CPU、耗電——先測而後優，勿憑感覺。
\end{enumerate}

\booksubsection{二、崩潰上報}

爾測試之時不崩，非謂用戶之機不崩也。上架之前，當接 Firebase Crashlytics（谷歌官方之物，免費）：配置數行，用戶崩潰則自動上報堆棧，爾於後台即可見「何機型、何版本、崩於何行」。無崩潰監控而上架者，猶閉目而駕車也。

\begin{tcbpractice}{來動手：將記帳本「交付」一次}
(1) 為 parseAmount 諸邏輯補單元測試三則；(2) 以 Logcat 親捕自造之崩潰一處而修之；(3) 生 release 簽名之 APK 一枚，致諸友朋安裝之（此乃「作品」之儀也）；(4) 力有餘則註冊開發者之號，以記帳本真呈上架。前三條今夕可竟。

\end{tcbpractice}

\begin{tcbsummary}{本章小結}
\begin{itemize}
  \item 調試之三技：Logcat 觀日誌、斷點觀變量、Layout Inspector 觀界面
  \item 讀崩潰：自堆棧之下而上，覓己包名之首行
  \item 測試三層：單元（秒級，先學此）/儀器/手動
  \item 發布必簽名：keystore 與密碼宜善備份之，失之則 App「易主」矣
  \item 上架六步：註冊 素材 上傳 AAB 數據安全表 審核 上架
\end{itemize}

\end{tcbsummary}
